\paragraph{RQ1：构建与容量。}
如何在 GPU 上构建 JHQ？其成本、常驻容量与构建峰值分别是多少（\cref{sec:7-7}）？

\paragraph{RQ2：正确且高效的执行。}
哪些执行选择能在保持分数／并列规则的同时加快搜索（\cref{sec:7-3,sec:7-4,sec:7-6}）？

\paragraph{RQ3：预算质量与总成本。}
无标签校准能否保持留出集质量，并收回自身成本（\cref{sec:7-5}）？使用 GT 标签的 oracle 是评估对照。

\paragraph{RQ4：具有竞争力的工作区间。}
按主机到主机边界计时，在相同召回率下 JHQ 何时领先？本文用独立调优的 E1 前沿及明确的显存定义回答（\cref{sec:7-2}）。Int8 图方法、高召回区间的胜负反转和未覆盖的区间也属于答案的一部分。

第 3 章将这些问题转化为受控比较，分别考察构建了什么、执行是否正确、工作负载需要多少计算，以及系统在哪些区间具有实际价值。
